% =============================================================================
% preguntas/p6-ncapas-calidad/solucion.tex — P6 (issue #11)
% Integrante: Enrique A Gonzalez Moreira
% Fragmento puro: se incluye con \include desde SemArq1_TeamBales.tex.
% Sin preámbulo: no declara clase de documento, inicio de documento ni paquetes.
% Fuente de trabajo: solucion.md (misma carpeta).
% Estructura del issue #11: \chapter con el título oficial + 3 \section (una por
% apartado del enunciado) + Notas de diseño + Notas de investigación.
% Citas: páginas impresas de Richards (2015) y Martin (2017); Richards y Ford
% (2020) no conserva paginación al extraer el texto del PDF y se cita por
% capítulo; Cockburn (2005) y Palermo (2008) son web sin paginación.
% Claves usadas, todas de core.bib: richards2015, richardsford2020, martin2017,
% cockburn2005, palermo2008. referencias.bib no declara entradas nuevas.
% Comillas: se usan comillas tipográficas literales, no \enquote, para no
% depender de csquotes (no está en la lista de paquetes de preamble.tex, #13).
% Tabla de calificaciones: obligatoria por el issue #11; ver la nota sobre su
% procedencia en «Notas de diseño».
% =============================================================================

\chapter{P6: Evaluación de N-Capas frente a Atributos de Calidad: Open Layers y Sinkhole Anti-pattern}
\label{sec:p6}

% Grupo local: \emergencystretch evita líneas desbordadas sin afectar a otros capítulos.
\begingroup\emergencystretch=3em \hbadness=2000

\section{Capas abiertas, aislamiento de capas y mantenibilidad a largo plazo}
\label{sec:p6-aislamiento}

El enunciado parte de una crítica conocida: la arquitectura N-Capas puntúa bajo en agilidad
general, facilidad de despliegue y escalabilidad \parencite[8-9]{richards2015}. Un estudiante
sostiene que esas carencias pueden mitigarse con capas abiertas (\emph{open layers}) y técnicas
como el \emph{fast-lane reader}. La afirmación es \textbf{tácticamente correcta pero
estratégicamente insuficiente}. Correcta, porque ambos mecanismos son reales y están
documentados: las propias fuentes los presentan para aliviar costes puntuales de la N-Capas
\parencite[cap.~10]{richardsford2020}. Insuficiente, porque ninguna de las dos técnicas corrige
la causa estructural de esas calificaciones bajas —la naturaleza monolítica y el acoplamiento
entre componentes \parencite[8-9]{richards2015}— y porque ambas se pagan con el único mecanismo
que la N-Capas ofrecía para gestionar el cambio con seguridad: el Aislamiento de Capas.

En la N-Capas cada capa es \textbf{cerrada} o \textbf{abierta}. Una capa cerrada no puede
saltarse: toda solicitud debe atravesar la capa inmediatamente inferior para llegar a la
siguiente \parencite[3]{richards2015}; \parencite[cap.~10]{richardsford2020}. Las capas cerradas
son las que materializan el \textbf{Aislamiento de Capas} (\emph{layers of isolation}): los
cambios hechos en una capa no impactan a componentes de otras siempre que los contratos entre
capas se mantengan, y cada capa es independiente de las demás y conoce poco o nada del interior
de las otras \parencite[3-4]{richards2015}.

Abrir una capa invierte esa garantía. Si la presentación puede acceder directamente a la
persistencia, un cambio en la persistencia impacta a la vez a la capa de negocio y a la de
presentación, produciendo “a very tightly coupled application with layer interdependencies” y
una arquitectura “very brittle, as well as difficult and expensive to change”
\parencite[cap.~10]{richardsford2020}; el mismo argumento, con idéntica conclusión de que el
sistema se vuelve “very hard and expensive to change”, está en \textcite[4]{richards2015}. El
\emph{fast-lane reader} es exactamente ese mecanismo: declarar abiertas las capas de negocio y
persistencia para que la presentación lea directamente la base de datos en consultas simples de
recuperación \parencite[cap.~10]{richardsford2020}.

A largo plazo, la mantenibilidad se juega en la gobernanza de esa decisión: “Failure to document
or properly communicate which layers in the architecture are open and closed (and why) usually
results in tightly coupled and brittle architectures that are very difficult to test, maintain,
and deploy” \parencite[5]{richards2015}. Con capas abiertas, el aislamiento deja de ser una
garantía de la arquitectura y pasa a depender de la disciplina del equipo: cada atajo autorizado
multiplica los caminos posibles de cada solicitud, de modo que conocer el impacto de un cambio
exige rastrear todos los flujos que cruzan capas. No sorprende que las mismas fuentes adviertan
que, a medida que crecen las aplicaciones en capas, “characteristics like maintainability,
agility, testability, and deployability are adversely affected”
\parencite[cap.~10]{richardsford2020}. La cuenta, en síntesis: la capa abierta compra
rendimiento local con la moneda de la mantenibilidad, justo el atributo que se quiere conservar
a largo plazo.

\section{Capas abiertas, sinkhole anti-pattern y la regla del 80/20}
\label{sec:p6-sinkhole}

El \textbf{architecture sinkhole anti-pattern} ocurre cuando las solicitudes atraviesan las
capas como \emph{pass-through} simple, sin lógica de negocio en ninguna de ellas
\parencite[7]{richards2015}; \parencite[cap.~10]{richardsford2020}. El ejemplo canónico: la
presentación pide datos básicos de un cliente; negocio y persistencia no hacen más que reenviar
la petición hasta una consulta SQL trivial; los datos suben por la pila sin agregarse,
calcularse ni transformarse. El coste no es inocuo: “unnecessary object instantiation and
processing, impacting both memory consumption and performance”
\parencite[cap.~10]{richardsford2020}.

Conviene precisar la dirección causal, porque las fuentes no dicen que las capas abiertas
\emph{causen} el sinkhole: dicen que son la \textbf{reacción} propuesta cuando el sinkhole
domina el tráfico. \textcite[7]{richards2015} sugiere considerar capas abiertas cuando la
proporción de peticiones \emph{pass-through} se invierte, advirtiendo de inmediato que “it will
be more difficult to control change due to the lack of layer isolation”;
\textcite[cap.~10]{richardsford2020} presentan abrir todas las capas como “another approach to
solving the architecture sinkhole anti-pattern”, con el mismo \emph{trade-off} de “increased
difficulty in managing change within the architecture”.

¿Por qué esa «solución» conduce de regreso al antipatrón? Porque no elimina su causa —capas
intermedias vacías de lógica— sino que \textbf{institucionaliza el rodeo}: la capa hueca sigue
existiendo, ahora como \emph{pass-through} permanente para las peticiones que aún la atraviesan,
y el atajo sancionado elimina el incentivo para escribir en ella lógica de negocio. Cuantas más
peticiones la salten, más se vacía; el sinkhole deja de ser un escenario puntual y se vuelve el
estado crónico de la arquitectura. El resultado combina lo peor de ambos mundos: el coste de
rendimiento del sinkhole y el acoplamiento de las capas abiertas.

La \textbf{regla del 80/20} desarma la propuesta del estudiante porque no es una perilla de
optimización sino un \textbf{diagnóstico de encaje del estilo}. Es típico que en torno al
20\,\% de las solicitudes sean \emph{pass-through} y el 80\,\% porten lógica de negocio; el
problema aparece cuando la proporción se invierte: “if 80 percent of the requests are sinkholes,
it [is] a good indicator that the layered architecture is not the correct architecture style for
the problem domain” \parencite[cap.~10]{richardsford2020}; \parencite[7]{richards2015}. Es
decir: si un sistema necesita capas abiertas y \emph{fast-lanes} por doquier para ser viable, la
regla no está autorizando la optimización; está dictaminando que la N-Capas no es la
arquitectura correcta para ese dominio. La respuesta razonable al diagnóstico es cambiar de
estilo, no optimizar el estilo equivocado, lo que conduce al tercer apartado.

\section{Hexagonal y Limpia frente a N-Capas para alta escalabilidad y agilidad}
\label{sec:p6-hexagonal-clean}

Las calificaciones bajas de la N-Capas son estructurales, no accidentales. La agilidad general
es baja porque hacer cambios sigue siendo “cumbersome and time-consuming [\ldots] because of the
monolithic nature of most implementations as well as the tight coupling of components”
\parencite[8]{richards2015}. El despliegue es bajo porque “One small change to a component can
require a redeployment of the entire application”, incompatible con la entrega continua
\parencite[8]{richards2015}; \textcite[cap.~10]{richardsford2020} puntúan la desplegabilidad
como “very low” —un cambio de tres líneas obliga a redesplegar la unidad completa— y la
elasticidad y escalabilidad como “very low (one star) [\ldots] primarily due to monolithic
deployments and the lack of architectural modularity”. La escalabilidad es baja por “the trend
toward tightly coupled and monolithic implementations” \parencite[9]{richards2015}.

\begin{table}[htbp]
  \centering
  \footnotesize
  \renewcommand{\arraystretch}{1.2}
  \begin{tabular}{@{}p{0.20\linewidth}p{0.24\linewidth}p{0.24\linewidth}p{0.24\linewidth}@{}}
    \toprule
    \textbf{Característica} & \textbf{N-Capas} & \textbf{Hexagonal} & \textbf{Limpia} \\
    \midrule
    \raggedright Agilidad general
      & \raggedright Baja: cambios lentos y engorrosos por la naturaleza monolítica y el
        acoplamiento \parencite[8]{richards2015}
      & \raggedright Alta: el núcleo se desarrolla y prueba aislado de dispositivos y bases de
        datos \parencite{cockburn2005}
      & \raggedright Alta: las reglas de negocio se prueban sin interfaz, base de datos ni
        servidor \parencite[202]{martin2017}
      \tabularnewline
    \raggedright Facilidad de despliegue
      & \raggedright Muy baja: un cambio de tres líneas obliga a redesplegar la unidad completa
        \parencite[cap.~10]{richardsford2020}
      & \raggedright Alta: cada adaptador es sustituible sin tocar el núcleo
        \parencite{cockburn2005}
      & \raggedright Alta: los detalles son \emph{plugins} reemplazables con mínimo esfuerzo
        \parencite[209]{martin2017}
      \tabularnewline
    \raggedright Escalabilidad
      & \raggedright Muy baja: despliegue monolítico y falta de modularidad arquitectónica
        \parencite[9]{richards2015}; \parencite[cap.~10]{richardsford2020}
      & \raggedright Alta: la partición por dominio permite escalar una capacidad sin arrastrar
        la pila \parencite{palermo2008}
      & \raggedright Alta: el núcleo no depende de la infraestructura que se escala
        \parencite[203]{martin2017}
      \tabularnewline
    \bottomrule
  \end{tabular}
  \caption{Atributos de calidad: N-Capas frente a Hexagonal y Limpia. Las valoraciones de la
    columna N-Capas son las puntuaciones publicadas en \textcite[8-9]{richards2015} y
    \textcite[cap.~10]{richardsford2020}. Hexagonal y Limpia no están puntuadas en esas fuentes:
    sus celdas son valoración propia, derivada de la Regla de Dependencia y de las propiedades
    que las fuentes citadas en cada celda atribuyen a esas arquitecturas.}
  \label{tab:p6-caracteristicas}
\end{table}

Las optimizaciones del estudiante no tocan ninguna de esas tres causas. Actúan sobre el
\textbf{camino de las solicitudes}, no sobre la \textbf{dirección de las dependencias} —el
negocio sigue dependiendo de la persistencia—, ni sobre la \textbf{unidad monolítica de
despliegue} —un cambio sigue redesplegando todo—, ni sobre el \textbf{tipo de partición}: la
N-Capas es una arquitectura técnicamente particionada, en la que el dominio (por ejemplo,
\texttt{customer}) está disperso por todas las capas, de modo que “a domain-driven design
approach does not work as well with the layered architecture style”
\parencite[cap.~10]{richardsford2020}. Peor: abrir capas añade dependencias transversales, es
decir, incrementa exactamente el acoplamiento que motiva las calificaciones bajas. El propio
libro que documenta estas técnicas lo dice del rendimiento: puede paliarse “through caching,
multithreading, and the like, but it is not a natural characteristic of this architecture style”
\parencite[cap.~10]{richardsford2020}. Las capas abiertas son lo mismo: un parche que no es
natural del estilo.

Las arquitecturas de dominio atacan la raíz, porque \textbf{invierten la dirección de las
dependencias}. La Regla de Dependencia de la Arquitectura Limpia es explícita: “Source code
dependencies must point only inward, toward higher-level policies. Nothing in an inner circle
can know anything at all about something in an outer circle” \parencite[203]{martin2017}, una
regla que se aplica siempre, con independencia del sentido del flujo de control, y que cruza
cada límite mediante el DIP \parencite[206]{martin2017}. En la Hexagonal de Cockburn, el núcleo
define \emph{puertos} y los \emph{adaptadores} externos los implementan, con el objetivo
explícito de que la aplicación se desarrolle y pruebe “in isolation from its eventual run-time
devices and databases”, permaneciendo “blissfully ignorant of the nature of the input device”
\parencite{cockburn2005}. En la Onion de Palermo, el Modelo de Dominio ocupa el centro, “only
coupled to itself”, con las interfaces —por ejemplo, las de repositorio— en el núcleo y el
comportamiento que toca bases de datos fuera \parencite{palermo2008}.

De esa única decisión estructural se derivan los atributos exigidos:

\begin{itemize}
  \item \textbf{Agilidad.} Las reglas de negocio se prueban “without the UI, database, web
        server, or any other external element” \parencite[202]{martin2017}: ciclos de prueba
        cortos y deterministas. Y los cambios de detalle no se propagan: al quedar obsoleto un
        elemento externo —base de datos, framework web—, “you can replace those obsolete
        elements with a minimum of fuss” \parencite[209]{martin2017}. En N-Capas, en cambio, un
        cambio menor en persistencia golpea a negocio y a presentación
        \parencite[4]{richards2015}.
  \item \textbf{Desplegabilidad.} Los detalles son \emph{plugins} del núcleo: sustituir o
        desplegar un adaptador no obliga a tocar el dominio; en N-Capas, el cambio de tres
        líneas redespliega la unidad entera \parencite[cap.~10]{richardsford2020};
        \parencite[8]{richards2015}.
  \item \textbf{Escalabilidad.} Al particionarse por dominio y no por capas técnicas, cada
        capacidad de negocio puede evolucionar, probarse, desplegarse y escalarse sin arrastrar
        la pila completa; en N-Capas, escalar una funcionalidad caliente equivale a escalar el
        monolito entero \parencite[9]{richards2015}.
\end{itemize}

Y la palabra «inherentemente» del enunciado es exacta: la diferencia entre ambos estilos no es
de grado sino de \textbf{dirección de las dependencias}, una propiedad definitoria y no un
parámetro ajustable. Abrir una capa no invierte dependencia alguna: solo autoriza más aristas en
la misma dirección hacia afuera. Para que una N-Capas cumpliera la Regla de Dependencia habría
que mover las interfaces al núcleo y hacer que la infraestructura las implemente, es decir,
reconstruirla como Hexagonal o Limpia. Por eso, ante requisitos de alta escalabilidad y
agilidad, las arquitecturas de dominio no son «una optimización más» de la N-Capas: son la
corrección de la causa que ninguna optimización de la N-Capas puede alcanzar.

\subsection*{Conclusión}

El estudiante acierta en que las técnicas existen y están documentadas, y falla en la
conclusión. Las capas abiertas y el \emph{fast-lane reader} compran rendimiento local con la
única moneda que la N-Capas tenía para el cambio controlado —el aislamiento—, institucionalizan
el sinkhole que pretendían evitar (la capa vaciada queda como \emph{pass-through} permanente y
sin incentivo para llenarse de lógica), y dejan intactas la naturaleza monolítica y el
acoplamiento que causan las calificaciones bajas en agilidad, despliegue y escalabilidad. La
regla del 80/20 es el veredicto: cuando esas técnicas dejan de ser excepción y se vuelven
necesidad, el estilo en capas no es el adecuado. Lo que exigen la alta escalabilidad y la
agilidad es estructural —dependencias apuntando hacia el dominio, detalles como \emph{plugins}
reemplazables, núcleo probable en aislamiento— y eso es, por definición, lo que ofrecen la
Hexagonal y la Limpia, y lo que la N-Capas no puede ofrecer «incluso con optimizaciones».

\section*{Notas de diseño}

\begin{itemize}
  \item \textbf{Dirección causal del sinkhole.} El enunciado sugiere que las capas abiertas
        «llevan» al sinkhole; las fuentes presentan la relación inversa: las capas abiertas se
        proponen como reacción al sinkhole, con un \emph{trade-off} explícito
        \parencite[7]{richards2015}; \parencite[cap.~10]{richardsford2020}. Para no contradecir
        el enunciado, la respuesta expone el bucle completo —capas vacías, sinkhole, apertura
        sancionada del rodeo, vaciado acelerado y aislamiento perdido—, que es la lectura fiel a
        las fuentes y responde a la preocupación real del apartado: que la «solución» degenera
        la arquitectura.
  \item \textbf{Procedencia de la tabla~\ref{tab:p6-caracteristicas}.} El issue~\#11 exige una
        tabla \texttt{booktabs} real con las puntuaciones de agilidad, desplegabilidad y
        escalabilidad comparando N-Capas con Hexagonal y Limpia.
        \textcite{richardsford2020} solo publican esas puntuaciones para los estilos que
        catalogan: no puntúan la Hexagonal ni la Limpia. Atribuirles valores habría sido una
        cita falsa, así que la columna de N-Capas reproduce las puntuaciones publicadas y las de
        Hexagonal y Limpia se declaran, en el pie de la tabla, como valoración propia derivada
        de las propiedades citadas celda a celda.
  \item \textbf{Estrategia de citas de Richards y Ford (2020).} El PDF disponible en
        \texttt{Bibliografias/} no conserva la paginación impresa al extraer el texto, así que
        ese libro se cita por capítulo y sección (“Layers of Isolation”, “Other Considerations”,
        “Architecture Characteristics Ratings”). \textcite{richards2015} y
        \textcite{martin2017} sí se citan con página exacta: las de Richards se verificaron
        contra los pies de página del propio reporte (numeración impresa del capítulo~1; en el
        PDF local aparece desplazada por las páginas preliminares) y las de Martin se
        contrastaron con las ya usadas en P4.
  \item \textbf{Doble cita Richards (2015) y Richards y Ford (2020).} El seminario atribuye las
        calificaciones a \textcite{richards2015}; \textcite{richardsford2020} es el tratamiento
        ampliado del mismo autor principal. Se citó el primero para las calificaciones del
        patrón y el segundo para el mecanismo de capas abiertas, el \emph{fast-lane reader} y la
        partición técnica, por ser donde el tratamiento es más completo.
  \item \textbf{DDD sin Evans.} El enunciado menciona «principios de DDD». Se sustentó con la
        partición por dominio \parencite[cap.~10]{richardsford2020} y el dominio al centro
        \parencite{palermo2008}; \parencite[203]{martin2017}, sin recurrir a Evans, que no
        figura entre las fuentes consultadas para esta pregunta.
  \item \textbf{Citas en inglés.} Se conservaron textuales en el idioma original, con la misma
        convención del resto del documento. Donde el original contiene una errata (“it a good
        indicator”), se corrigió con inserción editorial marcada: “it [is] a good indicator”.
\end{itemize}

\section*{Notas de investigación}

% =============================================================================
% investigacion/p6-open-layers/notas.tex — INV-P6 (issue #12)
% Integrante: Enrique A Gonzalez Moreira
% Fragmento puro: se incluye desde preguntas/p6-ncapas-calidad/solucion.tex
% mediante % =============================================================================
% investigacion/p6-open-layers/notas.tex — INV-P6 (issue #12)
% Integrante: Enrique A Gonzalez Moreira
% Fragmento puro: se incluye desde preguntas/p6-ncapas-calidad/solucion.tex
% mediante % =============================================================================
% investigacion/p6-open-layers/notas.tex — INV-P6 (issue #12)
% Integrante: Enrique A Gonzalez Moreira
% Fragmento puro: se incluye desde preguntas/p6-ncapas-calidad/solucion.tex
% mediante \input{investigacion/p6-open-layers/notas}.
% Estructura del issue #12: \section*{Notas de investigación: …} + Síntesis +
% Conceptos clave (tabla booktabs) + Puntos de conexión + Notas de diseño.
% Fuente de trabajo: notas.md (misma carpeta).
% Citas: páginas impresas de Richards (2015) y Martin (2017); Richards y Ford
% (2020) no conserva paginación al extraer el texto del PDF y se cita por
% capítulo; Cockburn (2005) y Palermo (2008) son web sin paginación.
% Todas las claves están en core.bib (richards2015, richardsford2020,
% martin2017, cockburn2005, palermo2008): fuentes.bib no declara nada.
% Comillas: se usan comillas tipográficas literales, no \enquote, para no
% depender de csquotes (no está en la lista de paquetes de preamble.tex, #13).
% =============================================================================

% Grupo local: \emergencystretch evita líneas desbordadas sin afectar a otros capítulos.
\begingroup\emergencystretch=3em \hbadness=2000
\section*{Notas de investigación: Open Layers, Sinkhole Anti-pattern y atributos de calidad en N-Capas}

\section{Síntesis}

La arquitectura N-Capas organiza el sistema en capas horizontales en las que cada capa conoce
únicamente a la inmediatamente inferior. Una capa es \textbf{cerrada} cuando ninguna solicitud
puede saltársela —toda petición debe atravesarla para alcanzar la siguiente— y \textbf{abierta}
cuando sí puede evitarse, accediendo directamente a una capa más profunda
\parencite[3]{richards2015}. Las capas cerradas son las que materializan el \emph{aislamiento de
capas}: los cambios realizados en una capa no impactan a los componentes de las demás mientras
se respeten los contratos, de modo que cada capa es independiente y conoce poco o nada del
interior de las otras \parencite[3-4]{richards2015}.

Abrir una capa compra rendimiento a costa de ese aislamiento. El caso típico es el
\emph{fast-lane reader}: se declaran abiertas las capas de negocio y persistencia para que la
presentación lea directamente la base de datos en consultas simples de recuperación
\parencite[cap.~10]{richardsford2020}. El precio es un acoplamiento transversal —un cambio en
persistencia pasa a impactar a la vez a negocio y a presentación— y una arquitectura
descrita como frágil y cara de cambiar \parencite[4]{richards2015}. Por eso la decisión es
sobre todo de gobernanza: no documentar ni comunicar qué capas son abiertas y cuáles cerradas,
y por qué, produce arquitecturas acopladas y frágiles, difíciles de probar, mantener y desplegar
\parencite[5]{richards2015}.

El \emph{architecture sinkhole anti-pattern} aparece cuando las solicitudes atraviesan las capas
como simple \emph{pass-through}, sin que ninguna aporte lógica de negocio: el coste es
instanciación y procesamiento innecesarios, con impacto en memoria y rendimiento
\parencite[cap.~10]{richardsford2020}. Importa la dirección causal: las fuentes no presentan las
capas abiertas como causa del sinkhole, sino como reacción propuesta cuando el sinkhole domina
el tráfico, con el aviso explícito de que entonces será más difícil controlar el cambio por la
falta de aislamiento \parencite[7]{richards2015}. La \textbf{regla del 80/20} es el criterio de
diagnóstico: es normal que en torno al 20\,\% de las peticiones sean \emph{pass-through} y el
80\,\% porten lógica, y si la proporción se invierte eso indica que la arquitectura en capas no
es el estilo correcto para ese dominio \parencite[7]{richards2015}.

Las calificaciones bajas de la N-Capas en agilidad, facilidad de despliegue y escalabilidad no
son accidentales: proceden de su naturaleza monolítica y del acoplamiento entre componentes
\parencite[8-9]{richards2015}. Para requisitos de alta escalabilidad y agilidad, la Hexagonal y
la Limpia son estructuralmente más adecuadas porque invierten la dirección de las dependencias:
el núcleo declara los puertos y la infraestructura los implementa \parencite{cockburn2005}, con
dependencias que apuntan solo hacia adentro \parencite[203]{martin2017} y un dominio que no se
acopla más que a sí mismo \parencite{palermo2008}.

\section{Conceptos clave}

\begin{table}[htbp]
  \centering
  \footnotesize
  \renewcommand{\arraystretch}{1.2}
  \begin{tabular}{@{}p{0.22\linewidth}p{0.46\linewidth}p{0.24\linewidth}@{}}
    \toprule
    \textbf{Concepto} & \textbf{Definición} & \textbf{Fuente} \\
    \midrule
    \raggedright Capa cerrada
      & \raggedright Capa que no puede saltarse: toda solicitud debe atravesarla para llegar a
        la capa inferior
      & \raggedright \textcite[3]{richards2015}; \textcite[cap.~10]{richardsford2020}
        \tabularnewline
    \raggedright Capa abierta
      & \raggedright Capa que una solicitud puede evitar, accediendo directamente a una capa más
        profunda
      & \raggedright \textcite[3]{richards2015}; \textcite[cap.~10]{richardsford2020}
        \tabularnewline
    \raggedright Aislamiento de capas
      & \raggedright Los cambios en una capa no impactan a las demás mientras se mantengan los
        contratos entre capas
      & \raggedright \textcite[3-4]{richards2015}; \textcite[cap.~10]{richardsford2020}
        \tabularnewline
    \raggedright Gobernanza de las capas abiertas
      & \raggedright No documentar qué capas son abiertas o cerradas produce arquitecturas
        acopladas y frágiles, difíciles de probar, mantener y desplegar
      & \raggedright \textcite[5]{richards2015} \tabularnewline
    \raggedright \emph{Fast-lane reader}
      & \raggedright Declarar abiertas negocio y persistencia para que la presentación lea
        directamente la base de datos en consultas simples de alto volumen
      & \raggedright \textcite[cap.~10]{richardsford2020} \tabularnewline
    \raggedright \emph{Architecture sinkhole anti-pattern}
      & \raggedright Solicitudes que atraviesan las capas como \emph{pass-through}, sin lógica
        en ninguna, con coste de instanciación y procesamiento
      & \raggedright \textcite[7]{richards2015}; \textcite[cap.~10]{richardsford2020}
        \tabularnewline
    \raggedright Regla del 80/20
      & \raggedright Si en torno al 80\,\% de las peticiones son \emph{pass-through}, el estilo
        en capas no es el adecuado para ese dominio
      & \raggedright \textcite[7]{richards2015}; \textcite[cap.~10]{richardsford2020}
        \tabularnewline
    \raggedright Partición técnica frente a partición por dominio
      & \raggedright En N-Capas el dominio queda disperso entre capas técnicas, por lo que un
        enfoque dirigido por el dominio encaja mal
      & \raggedright \textcite[cap.~10]{richardsford2020} \tabularnewline
    \raggedright Calificaciones de la N-Capas
      & \raggedright Agilidad y facilidad de despliegue bajas por la naturaleza monolítica y el
        acoplamiento; escalabilidad baja por la misma razón
      & \raggedright \textcite[8-9]{richards2015}; \textcite[cap.~10]{richardsford2020}
        \tabularnewline
    \raggedright Regla de Dependencia
      & \raggedright Las dependencias del código fuente apuntan solo hacia adentro, hacia las
        políticas de mayor nivel
      & \raggedright \textcite[203]{martin2017} \tabularnewline
    \raggedright Testabilidad del núcleo
      & \raggedright Las reglas de negocio se prueban sin interfaz de usuario, base de datos ni
        servidor
      & \raggedright \textcite[202]{martin2017} \tabularnewline
    \raggedright Puertos y adaptadores; dominio al centro
      & \raggedright El núcleo declara los puertos y la infraestructura los implementa; el
        Modelo de Dominio no se acopla más que a sí mismo
      & \raggedright \textcite{cockburn2005}; \textcite{palermo2008} \tabularnewline
    \bottomrule
  \end{tabular}
  \caption{Conceptos clave de la investigación INV-P6.}
  \label{tab:inv-p6-conceptos}
\end{table}

\section{Puntos de conexión con las preguntas}

\begin{itemize}
  \item \textbf{P6, primer apartado:} capa cerrada frente a capa abierta y aislamiento de capas
        dan el mecanismo exacto por el que abrir una capa degrada la mantenibilidad
        \parencite[3-5]{richards2015}.
  \item \textbf{P6, segundo apartado:} el sinkhole y la regla del 80/20 permiten mostrar que la
        apertura de capas es reacción al antipatrón, no su causa, y que la regla es un
        diagnóstico de encaje del estilo, no una perilla de optimización
        \parencite[7]{richards2015}.
  \item \textbf{P6, tercer apartado:} las calificaciones bajas \parencite[8-9]{richards2015} y
        la partición técnica \parencite[cap.~10]{richardsford2020} se contraponen a la inversión
        de dependencias de la Hexagonal y la Limpia \parencite[203]{martin2017}.
  \item \textbf{P1 y P2:} el mismo contraste N-Capas frente a Hexagonal/Onion sostiene el
        análisis de \emph{trade-offs} de P1 y la comparación de simetría de P2.
\end{itemize}

\section*{Notas de diseño}

\begin{itemize}
  \item \textbf{Dirección causal del sinkhole.} El enunciado de P6 sugiere que las capas
        abiertas llevan al sinkhole; las fuentes dicen lo inverso: son la reacción propuesta
        cuando el sinkhole domina \parencite[7]{richards2015}. La investigación documenta el
        bucle completo —capas vacías, sinkhole, apertura sancionada del rodeo, vaciado
        acelerado— para que la respuesta pueda ser fiel a las fuentes sin contradecir el
        enunciado.
  \item \textbf{Paginación de Richards y Ford (2020).} El PDF de \texttt{Bibliografias/} no
        conserva la paginación impresa al extraer el texto, así que ese libro se cita por
        capítulo y sección. Las páginas de \textcite{richards2015} se verificaron contra la
        numeración impresa del capítulo~1, que en el PDF local aparece desplazada por las
        páginas preliminares.
  \item \textbf{Doble cita de Richards.} \textcite{richards2015} aporta las calificaciones del
        patrón y \textcite{richardsford2020} el tratamiento ampliado del mismo autor principal
        (capas abiertas, \emph{fast-lane reader}, partición técnica). Se citan juntos donde
        coinciden y por separado donde cada uno es más explícito.
  \item \textbf{De la Torre Llorente y otros (2010).} Figura en \texttt{fuentes.md} como
        contexto del seminario sobre N-Capas orientada al dominio, pero no aporta ninguna
        afirmación de estas notas, así que no se cita.
  \item \textbf{DDD sin Evans.} El enunciado menciona principios de DDD. Se sostienen con la
        partición por dominio \parencite[cap.~10]{richardsford2020} y el dominio al centro
        \parencite{palermo2008}, sin recurrir a Evans, que no figura en la bibliografía
        consultada para esta investigación.
\end{itemize}
\endgroup
.
% Estructura del issue #12: \section*{Notas de investigación: …} + Síntesis +
% Conceptos clave (tabla booktabs) + Puntos de conexión + Notas de diseño.
% Fuente de trabajo: notas.md (misma carpeta).
% Citas: páginas impresas de Richards (2015) y Martin (2017); Richards y Ford
% (2020) no conserva paginación al extraer el texto del PDF y se cita por
% capítulo; Cockburn (2005) y Palermo (2008) son web sin paginación.
% Todas las claves están en core.bib (richards2015, richardsford2020,
% martin2017, cockburn2005, palermo2008): fuentes.bib no declara nada.
% Comillas: se usan comillas tipográficas literales, no \enquote, para no
% depender de csquotes (no está en la lista de paquetes de preamble.tex, #13).
% =============================================================================

% Grupo local: \emergencystretch evita líneas desbordadas sin afectar a otros capítulos.
\begingroup\emergencystretch=3em \hbadness=2000
\section*{Notas de investigación: Open Layers, Sinkhole Anti-pattern y atributos de calidad en N-Capas}

\section{Síntesis}

La arquitectura N-Capas organiza el sistema en capas horizontales en las que cada capa conoce
únicamente a la inmediatamente inferior. Una capa es \textbf{cerrada} cuando ninguna solicitud
puede saltársela —toda petición debe atravesarla para alcanzar la siguiente— y \textbf{abierta}
cuando sí puede evitarse, accediendo directamente a una capa más profunda
\parencite[3]{richards2015}. Las capas cerradas son las que materializan el \emph{aislamiento de
capas}: los cambios realizados en una capa no impactan a los componentes de las demás mientras
se respeten los contratos, de modo que cada capa es independiente y conoce poco o nada del
interior de las otras \parencite[3-4]{richards2015}.

Abrir una capa compra rendimiento a costa de ese aislamiento. El caso típico es el
\emph{fast-lane reader}: se declaran abiertas las capas de negocio y persistencia para que la
presentación lea directamente la base de datos en consultas simples de recuperación
\parencite[cap.~10]{richardsford2020}. El precio es un acoplamiento transversal —un cambio en
persistencia pasa a impactar a la vez a negocio y a presentación— y una arquitectura
descrita como frágil y cara de cambiar \parencite[4]{richards2015}. Por eso la decisión es
sobre todo de gobernanza: no documentar ni comunicar qué capas son abiertas y cuáles cerradas,
y por qué, produce arquitecturas acopladas y frágiles, difíciles de probar, mantener y desplegar
\parencite[5]{richards2015}.

El \emph{architecture sinkhole anti-pattern} aparece cuando las solicitudes atraviesan las capas
como simple \emph{pass-through}, sin que ninguna aporte lógica de negocio: el coste es
instanciación y procesamiento innecesarios, con impacto en memoria y rendimiento
\parencite[cap.~10]{richardsford2020}. Importa la dirección causal: las fuentes no presentan las
capas abiertas como causa del sinkhole, sino como reacción propuesta cuando el sinkhole domina
el tráfico, con el aviso explícito de que entonces será más difícil controlar el cambio por la
falta de aislamiento \parencite[7]{richards2015}. La \textbf{regla del 80/20} es el criterio de
diagnóstico: es normal que en torno al 20\,\% de las peticiones sean \emph{pass-through} y el
80\,\% porten lógica, y si la proporción se invierte eso indica que la arquitectura en capas no
es el estilo correcto para ese dominio \parencite[7]{richards2015}.

Las calificaciones bajas de la N-Capas en agilidad, facilidad de despliegue y escalabilidad no
son accidentales: proceden de su naturaleza monolítica y del acoplamiento entre componentes
\parencite[8-9]{richards2015}. Para requisitos de alta escalabilidad y agilidad, la Hexagonal y
la Limpia son estructuralmente más adecuadas porque invierten la dirección de las dependencias:
el núcleo declara los puertos y la infraestructura los implementa \parencite{cockburn2005}, con
dependencias que apuntan solo hacia adentro \parencite[203]{martin2017} y un dominio que no se
acopla más que a sí mismo \parencite{palermo2008}.

\section{Conceptos clave}

\begin{table}[htbp]
  \centering
  \footnotesize
  \renewcommand{\arraystretch}{1.2}
  \begin{tabular}{@{}p{0.22\linewidth}p{0.46\linewidth}p{0.24\linewidth}@{}}
    \toprule
    \textbf{Concepto} & \textbf{Definición} & \textbf{Fuente} \\
    \midrule
    \raggedright Capa cerrada
      & \raggedright Capa que no puede saltarse: toda solicitud debe atravesarla para llegar a
        la capa inferior
      & \raggedright \textcite[3]{richards2015}; \textcite[cap.~10]{richardsford2020}
        \tabularnewline
    \raggedright Capa abierta
      & \raggedright Capa que una solicitud puede evitar, accediendo directamente a una capa más
        profunda
      & \raggedright \textcite[3]{richards2015}; \textcite[cap.~10]{richardsford2020}
        \tabularnewline
    \raggedright Aislamiento de capas
      & \raggedright Los cambios en una capa no impactan a las demás mientras se mantengan los
        contratos entre capas
      & \raggedright \textcite[3-4]{richards2015}; \textcite[cap.~10]{richardsford2020}
        \tabularnewline
    \raggedright Gobernanza de las capas abiertas
      & \raggedright No documentar qué capas son abiertas o cerradas produce arquitecturas
        acopladas y frágiles, difíciles de probar, mantener y desplegar
      & \raggedright \textcite[5]{richards2015} \tabularnewline
    \raggedright \emph{Fast-lane reader}
      & \raggedright Declarar abiertas negocio y persistencia para que la presentación lea
        directamente la base de datos en consultas simples de alto volumen
      & \raggedright \textcite[cap.~10]{richardsford2020} \tabularnewline
    \raggedright \emph{Architecture sinkhole anti-pattern}
      & \raggedright Solicitudes que atraviesan las capas como \emph{pass-through}, sin lógica
        en ninguna, con coste de instanciación y procesamiento
      & \raggedright \textcite[7]{richards2015}; \textcite[cap.~10]{richardsford2020}
        \tabularnewline
    \raggedright Regla del 80/20
      & \raggedright Si en torno al 80\,\% de las peticiones son \emph{pass-through}, el estilo
        en capas no es el adecuado para ese dominio
      & \raggedright \textcite[7]{richards2015}; \textcite[cap.~10]{richardsford2020}
        \tabularnewline
    \raggedright Partición técnica frente a partición por dominio
      & \raggedright En N-Capas el dominio queda disperso entre capas técnicas, por lo que un
        enfoque dirigido por el dominio encaja mal
      & \raggedright \textcite[cap.~10]{richardsford2020} \tabularnewline
    \raggedright Calificaciones de la N-Capas
      & \raggedright Agilidad y facilidad de despliegue bajas por la naturaleza monolítica y el
        acoplamiento; escalabilidad baja por la misma razón
      & \raggedright \textcite[8-9]{richards2015}; \textcite[cap.~10]{richardsford2020}
        \tabularnewline
    \raggedright Regla de Dependencia
      & \raggedright Las dependencias del código fuente apuntan solo hacia adentro, hacia las
        políticas de mayor nivel
      & \raggedright \textcite[203]{martin2017} \tabularnewline
    \raggedright Testabilidad del núcleo
      & \raggedright Las reglas de negocio se prueban sin interfaz de usuario, base de datos ni
        servidor
      & \raggedright \textcite[202]{martin2017} \tabularnewline
    \raggedright Puertos y adaptadores; dominio al centro
      & \raggedright El núcleo declara los puertos y la infraestructura los implementa; el
        Modelo de Dominio no se acopla más que a sí mismo
      & \raggedright \textcite{cockburn2005}; \textcite{palermo2008} \tabularnewline
    \bottomrule
  \end{tabular}
  \caption{Conceptos clave de la investigación INV-P6.}
  \label{tab:inv-p6-conceptos}
\end{table}

\section{Puntos de conexión con las preguntas}

\begin{itemize}
  \item \textbf{P6, primer apartado:} capa cerrada frente a capa abierta y aislamiento de capas
        dan el mecanismo exacto por el que abrir una capa degrada la mantenibilidad
        \parencite[3-5]{richards2015}.
  \item \textbf{P6, segundo apartado:} el sinkhole y la regla del 80/20 permiten mostrar que la
        apertura de capas es reacción al antipatrón, no su causa, y que la regla es un
        diagnóstico de encaje del estilo, no una perilla de optimización
        \parencite[7]{richards2015}.
  \item \textbf{P6, tercer apartado:} las calificaciones bajas \parencite[8-9]{richards2015} y
        la partición técnica \parencite[cap.~10]{richardsford2020} se contraponen a la inversión
        de dependencias de la Hexagonal y la Limpia \parencite[203]{martin2017}.
  \item \textbf{P1 y P2:} el mismo contraste N-Capas frente a Hexagonal/Onion sostiene el
        análisis de \emph{trade-offs} de P1 y la comparación de simetría de P2.
\end{itemize}

\section*{Notas de diseño}

\begin{itemize}
  \item \textbf{Dirección causal del sinkhole.} El enunciado de P6 sugiere que las capas
        abiertas llevan al sinkhole; las fuentes dicen lo inverso: son la reacción propuesta
        cuando el sinkhole domina \parencite[7]{richards2015}. La investigación documenta el
        bucle completo —capas vacías, sinkhole, apertura sancionada del rodeo, vaciado
        acelerado— para que la respuesta pueda ser fiel a las fuentes sin contradecir el
        enunciado.
  \item \textbf{Paginación de Richards y Ford (2020).} El PDF de \texttt{Bibliografias/} no
        conserva la paginación impresa al extraer el texto, así que ese libro se cita por
        capítulo y sección. Las páginas de \textcite{richards2015} se verificaron contra la
        numeración impresa del capítulo~1, que en el PDF local aparece desplazada por las
        páginas preliminares.
  \item \textbf{Doble cita de Richards.} \textcite{richards2015} aporta las calificaciones del
        patrón y \textcite{richardsford2020} el tratamiento ampliado del mismo autor principal
        (capas abiertas, \emph{fast-lane reader}, partición técnica). Se citan juntos donde
        coinciden y por separado donde cada uno es más explícito.
  \item \textbf{De la Torre Llorente y otros (2010).} Figura en \texttt{fuentes.md} como
        contexto del seminario sobre N-Capas orientada al dominio, pero no aporta ninguna
        afirmación de estas notas, así que no se cita.
  \item \textbf{DDD sin Evans.} El enunciado menciona principios de DDD. Se sostienen con la
        partición por dominio \parencite[cap.~10]{richardsford2020} y el dominio al centro
        \parencite{palermo2008}, sin recurrir a Evans, que no figura en la bibliografía
        consultada para esta investigación.
\end{itemize}
\endgroup
.
% Estructura del issue #12: \section*{Notas de investigación: …} + Síntesis +
% Conceptos clave (tabla booktabs) + Puntos de conexión + Notas de diseño.
% Fuente de trabajo: notas.md (misma carpeta).
% Citas: páginas impresas de Richards (2015) y Martin (2017); Richards y Ford
% (2020) no conserva paginación al extraer el texto del PDF y se cita por
% capítulo; Cockburn (2005) y Palermo (2008) son web sin paginación.
% Todas las claves están en core.bib (richards2015, richardsford2020,
% martin2017, cockburn2005, palermo2008): fuentes.bib no declara nada.
% Comillas: se usan comillas tipográficas literales, no \enquote, para no
% depender de csquotes (no está en la lista de paquetes de preamble.tex, #13).
% =============================================================================

% Grupo local: \emergencystretch evita líneas desbordadas sin afectar a otros capítulos.
\begingroup\emergencystretch=3em \hbadness=2000
\section*{Notas de investigación: Open Layers, Sinkhole Anti-pattern y atributos de calidad en N-Capas}

\section{Síntesis}

La arquitectura N-Capas organiza el sistema en capas horizontales en las que cada capa conoce
únicamente a la inmediatamente inferior. Una capa es \textbf{cerrada} cuando ninguna solicitud
puede saltársela —toda petición debe atravesarla para alcanzar la siguiente— y \textbf{abierta}
cuando sí puede evitarse, accediendo directamente a una capa más profunda
\parencite[3]{richards2015}. Las capas cerradas son las que materializan el \emph{aislamiento de
capas}: los cambios realizados en una capa no impactan a los componentes de las demás mientras
se respeten los contratos, de modo que cada capa es independiente y conoce poco o nada del
interior de las otras \parencite[3-4]{richards2015}.

Abrir una capa compra rendimiento a costa de ese aislamiento. El caso típico es el
\emph{fast-lane reader}: se declaran abiertas las capas de negocio y persistencia para que la
presentación lea directamente la base de datos en consultas simples de recuperación
\parencite[cap.~10]{richardsford2020}. El precio es un acoplamiento transversal —un cambio en
persistencia pasa a impactar a la vez a negocio y a presentación— y una arquitectura
descrita como frágil y cara de cambiar \parencite[4]{richards2015}. Por eso la decisión es
sobre todo de gobernanza: no documentar ni comunicar qué capas son abiertas y cuáles cerradas,
y por qué, produce arquitecturas acopladas y frágiles, difíciles de probar, mantener y desplegar
\parencite[5]{richards2015}.

El \emph{architecture sinkhole anti-pattern} aparece cuando las solicitudes atraviesan las capas
como simple \emph{pass-through}, sin que ninguna aporte lógica de negocio: el coste es
instanciación y procesamiento innecesarios, con impacto en memoria y rendimiento
\parencite[cap.~10]{richardsford2020}. Importa la dirección causal: las fuentes no presentan las
capas abiertas como causa del sinkhole, sino como reacción propuesta cuando el sinkhole domina
el tráfico, con el aviso explícito de que entonces será más difícil controlar el cambio por la
falta de aislamiento \parencite[7]{richards2015}. La \textbf{regla del 80/20} es el criterio de
diagnóstico: es normal que en torno al 20\,\% de las peticiones sean \emph{pass-through} y el
80\,\% porten lógica, y si la proporción se invierte eso indica que la arquitectura en capas no
es el estilo correcto para ese dominio \parencite[7]{richards2015}.

Las calificaciones bajas de la N-Capas en agilidad, facilidad de despliegue y escalabilidad no
son accidentales: proceden de su naturaleza monolítica y del acoplamiento entre componentes
\parencite[8-9]{richards2015}. Para requisitos de alta escalabilidad y agilidad, la Hexagonal y
la Limpia son estructuralmente más adecuadas porque invierten la dirección de las dependencias:
el núcleo declara los puertos y la infraestructura los implementa \parencite{cockburn2005}, con
dependencias que apuntan solo hacia adentro \parencite[203]{martin2017} y un dominio que no se
acopla más que a sí mismo \parencite{palermo2008}.

\section{Conceptos clave}

\begin{table}[htbp]
  \centering
  \footnotesize
  \renewcommand{\arraystretch}{1.2}
  \begin{tabular}{@{}p{0.22\linewidth}p{0.46\linewidth}p{0.24\linewidth}@{}}
    \toprule
    \textbf{Concepto} & \textbf{Definición} & \textbf{Fuente} \\
    \midrule
    \raggedright Capa cerrada
      & \raggedright Capa que no puede saltarse: toda solicitud debe atravesarla para llegar a
        la capa inferior
      & \raggedright \textcite[3]{richards2015}; \textcite[cap.~10]{richardsford2020}
        \tabularnewline
    \raggedright Capa abierta
      & \raggedright Capa que una solicitud puede evitar, accediendo directamente a una capa más
        profunda
      & \raggedright \textcite[3]{richards2015}; \textcite[cap.~10]{richardsford2020}
        \tabularnewline
    \raggedright Aislamiento de capas
      & \raggedright Los cambios en una capa no impactan a las demás mientras se mantengan los
        contratos entre capas
      & \raggedright \textcite[3-4]{richards2015}; \textcite[cap.~10]{richardsford2020}
        \tabularnewline
    \raggedright Gobernanza de las capas abiertas
      & \raggedright No documentar qué capas son abiertas o cerradas produce arquitecturas
        acopladas y frágiles, difíciles de probar, mantener y desplegar
      & \raggedright \textcite[5]{richards2015} \tabularnewline
    \raggedright \emph{Fast-lane reader}
      & \raggedright Declarar abiertas negocio y persistencia para que la presentación lea
        directamente la base de datos en consultas simples de alto volumen
      & \raggedright \textcite[cap.~10]{richardsford2020} \tabularnewline
    \raggedright \emph{Architecture sinkhole anti-pattern}
      & \raggedright Solicitudes que atraviesan las capas como \emph{pass-through}, sin lógica
        en ninguna, con coste de instanciación y procesamiento
      & \raggedright \textcite[7]{richards2015}; \textcite[cap.~10]{richardsford2020}
        \tabularnewline
    \raggedright Regla del 80/20
      & \raggedright Si en torno al 80\,\% de las peticiones son \emph{pass-through}, el estilo
        en capas no es el adecuado para ese dominio
      & \raggedright \textcite[7]{richards2015}; \textcite[cap.~10]{richardsford2020}
        \tabularnewline
    \raggedright Partición técnica frente a partición por dominio
      & \raggedright En N-Capas el dominio queda disperso entre capas técnicas, por lo que un
        enfoque dirigido por el dominio encaja mal
      & \raggedright \textcite[cap.~10]{richardsford2020} \tabularnewline
    \raggedright Calificaciones de la N-Capas
      & \raggedright Agilidad y facilidad de despliegue bajas por la naturaleza monolítica y el
        acoplamiento; escalabilidad baja por la misma razón
      & \raggedright \textcite[8-9]{richards2015}; \textcite[cap.~10]{richardsford2020}
        \tabularnewline
    \raggedright Regla de Dependencia
      & \raggedright Las dependencias del código fuente apuntan solo hacia adentro, hacia las
        políticas de mayor nivel
      & \raggedright \textcite[203]{martin2017} \tabularnewline
    \raggedright Testabilidad del núcleo
      & \raggedright Las reglas de negocio se prueban sin interfaz de usuario, base de datos ni
        servidor
      & \raggedright \textcite[202]{martin2017} \tabularnewline
    \raggedright Puertos y adaptadores; dominio al centro
      & \raggedright El núcleo declara los puertos y la infraestructura los implementa; el
        Modelo de Dominio no se acopla más que a sí mismo
      & \raggedright \textcite{cockburn2005}; \textcite{palermo2008} \tabularnewline
    \bottomrule
  \end{tabular}
  \caption{Conceptos clave de la investigación INV-P6.}
  \label{tab:inv-p6-conceptos}
\end{table}

\section{Puntos de conexión con las preguntas}

\begin{itemize}
  \item \textbf{P6, primer apartado:} capa cerrada frente a capa abierta y aislamiento de capas
        dan el mecanismo exacto por el que abrir una capa degrada la mantenibilidad
        \parencite[3-5]{richards2015}.
  \item \textbf{P6, segundo apartado:} el sinkhole y la regla del 80/20 permiten mostrar que la
        apertura de capas es reacción al antipatrón, no su causa, y que la regla es un
        diagnóstico de encaje del estilo, no una perilla de optimización
        \parencite[7]{richards2015}.
  \item \textbf{P6, tercer apartado:} las calificaciones bajas \parencite[8-9]{richards2015} y
        la partición técnica \parencite[cap.~10]{richardsford2020} se contraponen a la inversión
        de dependencias de la Hexagonal y la Limpia \parencite[203]{martin2017}.
  \item \textbf{P1 y P2:} el mismo contraste N-Capas frente a Hexagonal/Onion sostiene el
        análisis de \emph{trade-offs} de P1 y la comparación de simetría de P2.
\end{itemize}

\section*{Notas de diseño}

\begin{itemize}
  \item \textbf{Dirección causal del sinkhole.} El enunciado de P6 sugiere que las capas
        abiertas llevan al sinkhole; las fuentes dicen lo inverso: son la reacción propuesta
        cuando el sinkhole domina \parencite[7]{richards2015}. La investigación documenta el
        bucle completo —capas vacías, sinkhole, apertura sancionada del rodeo, vaciado
        acelerado— para que la respuesta pueda ser fiel a las fuentes sin contradecir el
        enunciado.
  \item \textbf{Paginación de Richards y Ford (2020).} El PDF de \texttt{Bibliografias/} no
        conserva la paginación impresa al extraer el texto, así que ese libro se cita por
        capítulo y sección. Las páginas de \textcite{richards2015} se verificaron contra la
        numeración impresa del capítulo~1, que en el PDF local aparece desplazada por las
        páginas preliminares.
  \item \textbf{Doble cita de Richards.} \textcite{richards2015} aporta las calificaciones del
        patrón y \textcite{richardsford2020} el tratamiento ampliado del mismo autor principal
        (capas abiertas, \emph{fast-lane reader}, partición técnica). Se citan juntos donde
        coinciden y por separado donde cada uno es más explícito.
  \item \textbf{De la Torre Llorente y otros (2010).} Figura en \texttt{fuentes.md} como
        contexto del seminario sobre N-Capas orientada al dominio, pero no aporta ninguna
        afirmación de estas notas, así que no se cita.
  \item \textbf{DDD sin Evans.} El enunciado menciona principios de DDD. Se sostienen con la
        partición por dominio \parencite[cap.~10]{richardsford2020} y el dominio al centro
        \parencite{palermo2008}, sin recurrir a Evans, que no figura en la bibliografía
        consultada para esta investigación.
\end{itemize}
\endgroup

\endgroup
